\chapter[INTRODUCTION]{Introduction}
\label{ch:introduction}

\section{Background}
\label{sec:background}
% -----------------------------------------------------------------
% Funnel, roughly 3-5 paragraphs:
%   1. Corrosion as an industrial and economic problem (asset
%      integrity, offshore/marine context, cost of failure).
%   2. How inspection is done now -- manual visual inspection, its
%      cost, subjectivity and access limitations.
%   3. Automated vision-based corrosion detection and where it has
%      got to; note that it is overwhelmingly RGB-only.
%   4. Why a single modality is physically insufficient: colour and
%      texture alone confuse corrosion with coating, staining, dirt
%      and shadow, and say nothing about material loss.
%   5. Land on: complementary sensing (thermal for subsurface,
%      LiDAR for geometry) as the direction this work takes.
% Do NOT argue the gap here -- that is the Problem Statement.
% -----------------------------------------------------------------

\note{Background to be drafted}


\section{Problem Statement}
\label{sec:problem-statement}
% -----------------------------------------------------------------
% This is the three-leg gap argument. Keep the legs explicit and
% separable -- examiners will attack them one at a time.
%
%   Leg 1 (domain gap): corrosion and marine/offshore inspection is
%       overwhelmingly RGB-only or dual-modality; trimodal is rare.
%   Leg 2 (sensor gap): the closest anchor work substitutes laser
%       profilometry for true LiDAR; scanning geometry, range and
%       field deployability differ materially.
%   Leg 3 (ablation gap): fusion is widely assumed to help, but few
%       primary studies test whether each added modality earns its
%       place. Where they do, gains are modest -- and that is a
%       finding, not a weakness in the proposal.
%
% The sharpened one-line version from the literature sweep:
%   the trimodal sensing is solved (calibration literature), the
%   bimodal corrosion detection is solved, and nobody has joined
%   the two.
% -----------------------------------------------------------------

\note{Problem statement to be drafted}


\section{Research Questions}
\label{sec:research-questions}
% -----------------------------------------------------------------
% Two or three questions, each one answerable by an objective below.
% Keep them interrogative and testable. Descriptive -> relational ->
% causal, per the CGS guidance deck.
% -----------------------------------------------------------------

\note{Research questions to be drafted}

% \begin{enumerate}
%   \item ...
%   \item ...
%   \item ...
% \end{enumerate}


\section{Research Aim and Objectives}
\label{sec:aim-objectives}
% -----------------------------------------------------------------
% One aim sentence, then exactly THREE objectives -- the template's
% methodology diagram maps three objectives onto five milestones,
% so three is the number the document expects.
% Each objective must start with a measurable verb (to develop, to
% evaluate, to quantify) and must map to a research question.
% -----------------------------------------------------------------

\subsection*{Aim}
\note{Aim statement to be drafted}

\subsection*{Objectives}
\begin{enumerate}[label=\textbf{Objective \arabic*:},leftmargin=*,align=left]
  \item \note{objective 1}
  \item \note{objective 2}
  \item \note{objective 3}
\end{enumerate}


\section{Significance and Contribution of Research}
\label{sec:significance}
% -----------------------------------------------------------------
% Separate the two:
%   Contribution -- what is new (the join that has not been made,
%       the ablation protocol, any dataset produced).
%   Significance -- who benefits and how (asset integrity practice,
%       inspection cost, the fusion literature's evidence base).
% Be honest about the negative-result case: if the third modality
% does not earn its place, that is still a publishable contribution.
% -----------------------------------------------------------------

\note{Significance and contribution to be drafted}


\section{Scope of Study}
\label{sec:scope}
% -----------------------------------------------------------------
% State the boundaries plainly -- this section is where examiners
% look for realism about a two-year MSc.
% In scope: the three modalities; the ablation protocol.
% Explicitly OUT of scope: real-time adaptive learning.
% Still open, decide before the defence: industry/application domain
% (offshore is the candidate), detection vs segmentation, whether
% severity quantification is attempted, and hyperspectral -- the
% sweep found a substantial hyperspectral corrosion literature, so
% say why thermal was chosen over it or the choice looks unexamined.
% -----------------------------------------------------------------

\note{Scope to be drafted}
